\section{Lecture 23 (December 4th)}
\begin{recall}
(Euler's constant) We have defined $f(z)=\prod^{\infty }_{n=1}(1+z/n)e^{-z/n}$ and $f(z-1)=e^{g(z)}zf(z)$. As $g'(z)=0$, we set $g(z)=\gamma $ and find:
\[g(z)=\gamma =\lim _{n\rightarrow \infty }\Big(1+\dfrac{1}{2}+\ldots +\dfrac{1}{n}-\ln(n+1)\Big)\]
\end{recall}
\vspace{2ex}
\begin{defi}
(Weierstrass gamma function) The Weierstrass gamma function is defined as:
\[\mathrm{\Gamma}(z)=\dfrac{1}{e^{\gamma z }zf(z)}=\dfrac{e^{-\gamma z}}{z}\prod^{\infty }_{n=1}\Big(1+\dfrac{z}{n}\Big)^{-1}e^{z/n}\]
where $e^{\gamma }zf(z)$ is a entire function with simple zeros at $0,-1,-2,\ldots $. Therefore, the function has simple poles at these points. 
\[\dfrac{1}{\mathrm{\Gamma} (z)}=e^{\gamma z}zf(z)=e^{\gamma z}z\prod^{\infty }_{n=1}\Big(1+\dfrac{z}{n}\Big)e^{-z/n}\]
From $f(z-1)=e^{\gamma }zf(z)$,
\[f(z)=e^{\gamma }(z+1)f(z+1)\]
so that 
\[\mathrm{\Gamma} (z+1)=\dfrac{1}{e^{\gamma (z+1)}(z+1)f(z+1)}=\dfrac{1}{e^{\gamma z}f(z)}=\dfrac{z}{e^{\gamma z}zf(z)}=z\mathrm{\Gamma} (z)\]
This tells us that $\mathrm{\Gamma} (z+1)=z\mathrm{\Gamma} (z)$ for all $z\in {\bm C}\backslash \{0,-1,-2,\ldots \}$. Meanwhile,
\[\mathrm{\Gamma} (1)=\dfrac{1}{e^{\gamma }f(1)}=\dfrac{1}{e^{\gamma }e^{-\gamma }}=1\]
and $\mathrm{\Gamma} (n+1)=n!$.
\end{defi}
\vspace{2ex}
\begin{defi}
(Euler Gamma function) Now, the Euler Gamma function is defined for $x\in {\bm R}$ and $x>0$ as:
\[\mathrm{\Gamma} (x)=\int ^{\infty }_{0}t^{x-1}e^{-t}\,dt\]
If $x\leq 0$,
\[\int ^{\infty }_{0}t^{x-1}e^{-t}\,dt\geq \int ^{1}_{0}t^{x-1}e^{-t}\,dt\geq \int ^{1}_{0}t^{x-1}e^{-1}\,dt\geq \int ^{1}_{0}\dfrac{1}{t}\dfrac{1}{e}\,dt=\infty \]
Meanwhile, for $z\in {\bm C}$ with $\mathrm{Re}(z)>0$,
\[\mathrm{\Gamma} (z)=\int ^{\infty }_{0}t^{z-1}e^{-t}\,dt\]
and with $z=x+iy$,
\[|t^{z}|=|e^{z\ln t}|=e^{x\ln t}=t^{x}\]
tells us that the integral is well-defined. The easiest way to see that the Euler Gamma function is a analytic continuation of the Weierstrass gamma function, we use Euler's identity. 
\end{defi}
\vspace{2ex}
\begin{recall}
Recall that
\[\dfrac{\sin \pi z}{\pi z}=\prod^{\infty }_{n=1}\Big(1-\dfrac{z^2}{n^2}\Big)=f(z)f(-z)\]
so that 
\[\sin \pi z=\pi z f(z)f(-z)=\pi z\dfrac{1}{e^{\gamma z}z\mathrm{\Gamma} (z)}\dfrac{1}{e^{-\gamma z}(-z)\mathrm{\Gamma} (-z)}=\dfrac{\pi }{\mathrm{\Gamma} (z)(-z)\mathrm{\Gamma} (-z)}=\dfrac{\pi }{\mathrm{\Gamma} (z)\mathrm{\Gamma} (1-z)}\]
That is,
\[\mathrm{\Gamma} (z)\mathrm{\Gamma} (1-z)=\dfrac{\pi }{\sin \pi z}\]
\end{recall}
\vspace{2ex}
\begin{thm}
(Euler identity) We see that
\[\mathrm{\Gamma} (z)\mathrm{\Gamma} (1-z)=\dfrac{\pi }{\sin \pi z} \]
for $z\in {\bm Z}$.
\end{thm}
\vspace{2ex}
\begin{proof}
For $0<s<1$,
\[\mathrm{\Gamma} (s)=\int ^{\infty }_{0}t^{s-1}e^{-t}\,dt\qquad\&\qquad \mathrm{\Gamma} (1-s)=\int ^{\infty }_{0}u^{-s}e^{-u}\,du\]
Putting $u=ty$ for $t>0$,
\[\mathrm{\Gamma} (1-s)=t\int ^{\infty }_{0}(ty)^{-s}e^{-ty}\,dy\]
Consider 
\[\int ^{\infty }_{0}t^{s-1}e^{-t}\Big[\int ^{\infty }_{0}t(ty)^{-s}e^{-ty}\,dy\Big]\,dt=\mathrm{\Gamma} (s)\mathrm{\Gamma} (1-s)\]
which is also equal to
\[\int ^{\infty }_{0}\int ^{\infty }_{0}e^{-(1+y)}y^{-s}\,dt\,dy=\int ^{\infty }_{0}\dfrac{y^{-s}}{1+y}\,dy=\dfrac{\pi }{\sin [(1-s) \pi ]}=\dfrac{\pi }{\sin 	\pi s}\]
Therefore,
\[\mathrm{\Gamma} (s)\mathrm{\Gamma} (1-s)=\dfrac{\pi }{\sin \pi s}\]
At this stage, using the identity theorem along the real line,
\[\mathrm{\Gamma} (z)\mathrm{\Gamma} (1-z)=\dfrac{\pi }{\sin \pi z}\]
for $0<\mathrm{Re}(z)<1$. 
\end{proof}
\vspace{2ex}
\begin{defi}
(Euler zeta function) The Euler zeta function is defined as:
\[\xi (s)=\sum ^{\infty }_{n=1}\dfrac{1}{n^{s}}\]
for $s>1$ if real or $\mathrm{Res}(s)<1$ if complex. Consider $\mathrm{Re}(s)=\sigma >1$.
\[\begin{align*}
\int ^{\infty }_{0}\dfrac{x^{s-1}}{e^{x}-1}\,dx=&\int ^{\infty }_{0}x^{s-1}\sum ^{\infty }_{n=1}e^{-nx}\,dx=\sum ^{\infty }_{n=1}\int ^{\infty }_{0}x^{s-1}e^{-nx}\,dx\\=&\sum ^{\infty }_{n=1}\int ^{\infty }_{0}x^{s-1}e^{-nx}\,dx
=\sum ^{\infty }_{n=1}\dfrac{1}{n^{s}}\int ^{\infty }_{0}y^{s-1}e^{-y}\,dy\\
=&\xi (s)\mathrm{\Gamma} (s)
\end{align*}\]
by setting $nx=y$. Thus we have the important identity that
\[\xi (s)=\dfrac{1}{\mathrm{\Gamma} (s)}\int ^{\infty }_{0}\dfrac{x^{s-1}}{e^{x}-1}\,dx\]
\end{defi}
\vspace{2ex}
\begin{defi}
(Dirichlet eta function) The function is defined as:
\[\eta (s)=\sum ^{\infty }_{n=1}(-1)^{n+1}\dfrac{1}{n^{s}}=1-\dfrac{1}{2^{s}}+\dfrac{1}{3^{s}}+\ldots \]
which converges in $\mathrm{Re}(s)>0$. This is closely related to $\xi (s)$ as 
\[\xi (s)-\eta (s)=2\sum ^{\infty }_{n=1}\dfrac{1}{(2n)^{s}}=2^{1-s}\xi (s)\]
and 
\[\xi (s)=(1-2^{1-s})^{-1}\eta (s)\]
when $\mathrm{Re}(s)>1$. Notice that this version of the zeta function has no roots in $0<\mathrm{Re}(s)<1$ unless
\[s=\dfrac{1}{2}+it\]
\end{defi}
\vspace{2ex}
\begin{lem}
(Prime number theorem) The prime number theorem states that:
\[\lim _{x\rightarrow \infty }\dfrac{\pi (x)\log x}{x}=1\]
\end{lem}
\vspace{2ex}

